% Einheit 2 von 3 – Bedingungen mit Ableitung (bank/rekonstruktion-von-funktionsgleichungen/e2.jsonl, zusammenbau v0.1)
%% TODO Zeitmarke Einheit 2: Marken-Zeile nicht lesbar
%% TODO Zweigzeile Teil 1: was hier gelernt wird (Satz in Schülersprache aus dem Einheitstitel) – Einheit 2
\einheitenkopf[e2]{Einheit 2 von 3 · Bedingungen mit Ableitung}
\zweigzeile{Abitur GK}
\setcounter{aufgabe}{10}

%% TODO Titel: Ich-kann-Satz fehlt in der Bank (Platzhalter: Kettenname) – Nr. 11
%% TODO Anweisung: ein Satz über der ersten Teilaufgabe fehlt in der Bank – Nr. 11
\begin{aufgabe}{Bedingungen mit Ableitung}
%% TODO \rechenplatz{n}: Schrittzahl des Grundfalls fehlt in der Bank (nur bei fester Schreibform, 2.3 b)
\begin{teile}
\teil „Der Graph von $f$ hat den Hochpunkt $H(3|1)$.“ Wie viele Gleichungen stecken in diesem Satz? \\ \kreuz{eine Gleichung} \\ \kreuz{zwei Gleichungen}
\teil Der Graph einer quadratischen Funktion $f$ verläuft durch $A(0|-1)$ und hat den Tiefpunkt $T(1|-3)$. Bestimme eine Funktionsgleichung von $f$. f(x) = \leerfeld
\teil Die Höhe einer Sonnenblume wird in den ersten Wochen durch $h(t) = a \cdot t^3 + b \cdot t^2$ beschrieben ($t$ in Wochen, $h(t)$ in cm). Nach 4 Wochen ist sie 64 cm hoch und wächst in diesem Moment mit 24 cm pro Woche. Bestimme eine Funktionsgleichung von $h$. \hfill (Abitur 2019 GK) \\ h(t) = \leerfeld
\teil Der Graph einer quadratischen Funktion $f$ verläuft durch den Ursprung; die Tangente an den Graphen im Punkt $(1|f(1))$ hat die Gleichung $y = 3x + 1$. Bestimme einen Funktionsterm von $f$. \hfill (Abitur 2020 GK) \\ f(x) = \leerfeld
\teil Bei einer Dirtbike-Schanze endet die Absprungrampe $r$ mit $r(x) = 0{,}1x^2 + 2{,}5$ ($x \le 0$, alle Angaben in m) im Punkt $S(0|2{,}5)$; dort geht sie ohne Knick in die Flugbahn über, die durch eine quadratische Funktion $f$ beschrieben wird. Der Landehügel wird durch $g(x) = 2{,}5 - 0{,}5x$ beschrieben. Bei $x = 4$ ist der Fahrer 1,2 m senkrecht über dem Landehügel. Bestimme eine Funktionsgleichung der Flugbahn $f$. \hfill (Abitur 2018 GK) \\ f(x) = \leerfeld
\teil Ein Rundbogen über einem Gartenweg hat 4 m Spannweite und wird durch eine zur $y$-Achse symmetrische Parabel $p$ beschrieben, deren Nullstellen die Fußpunkte sind (Angaben in m). An den Fußpunkten hat der Bogen die Steigungswinkel $45^\circ$ und $-45^\circ$. Bestimme eine Gleichung von $p$ und die Höhe des Bogens. p(x) = \leerfeld , Höhe \leerfeld
\teil Gegeben ist $f'(x) = 3x \cdot (2 - x)$. Die Gerade mit $t(x) = 3x + 2$ ist Tangente an den Graphen von $f$. Bestimme eine mögliche Funktionsgleichung von $f$. \hfill (Abitur 2023 LK) \\ f(x) = \leerfeld
\teil Das Wasservolumen eines Stausees nach dem Schließen der Schleuse wird durch $V(t) = 50 - 20 \cdot e^{k \cdot t}$ beschrieben ($t$ in Tagen, $V(t)$ in Millionen m³). Zu Beginn nimmt das Volumen um 4 Millionen m³ pro Tag zu. Bestimme $k$. \hfill (Abitur 2026 LK) \\ k = \leerfeld
\teil Über einen Bach führt eine Holzbrücke. Sie soll durch eine ganzrationale Funktion $p$ vierten Grades beschrieben werden, deren Graph symmetrisch zur $y$-Achse ist. Die Brücke hat eine Spannweite von 4 m, ist in der Mitte 1,8 m hoch und hat an den beiden Enden die Steigungswinkel $45^\circ$ bzw. $-45^\circ$. Bestimme eine Gleichung dieser Funktion. \hfill (Abitur 2018 LK) \\ p(x) = \leerfeld
\end{teile}
\end{aufgabe}

%% TODO Titel: Ich-kann-Satz fehlt in der Bank (Platzhalter: Kettenname) – Nr. 12
%% TODO Anweisung: ein Satz über der ersten Teilaufgabe fehlt in der Bank – Nr. 12
\begin{aufgabe}{Zwei Parameter einer Exponentialfunktion aus Funktionswert und Änderungsrate an einer Stelle bestimmen}
\begin{teile}
\teil Eine Suppe wird vom Herd genommen und kühlt ab. Ab $t = 5$ ($t$ in Minuten) wird ihre Temperatur durch $w(t) = 20 + b \cdot e^{c \cdot t}$ beschrieben ($w(t)$ in °C). Zu Beginn des Abkühlens beträgt die Temperatur 80 °C und die momentane Änderungsrate $-6$ °C pro Minute. Bestimme $b$ und $c$; runde $b$ auf zwei Stellen. \hfill (Abitur 2017 GK) \\ b ≈ \leerfeld , c = \leerfeld
\end{teile}
\end{aufgabe}

%% TODO Titel: Ich-kann-Satz fehlt in der Bank (Platzhalter: Kettenname) – Nr. 13
%% TODO Anweisung: ein Satz über der ersten Teilaufgabe fehlt in der Bank – Nr. 13
\begin{aufgabe}{Bedingungen mit Ableitung – Fehler finden, Begründen}
\begin{teile}
\teil Ich finde den Fehler: Gesucht ist eine quadratische Funktion $f$, deren Graph durch $A(0|-2)$ verläuft und den Hochpunkt $H(2|2)$ hat. Nina rechnet: \rechnung{c &= -2 \\ 4a + 2b - 2 &= 2 \\ b &= 0 && \text{(gewählt)} \\ a &= 1}
\teil Begründe, warum die Angabe „$H(3|5)$ ist Hochpunkt des Graphen von $f$“ zwei Gleichungen liefert.
\end{teile}
\end{aufgabe}
